% GENERATED FILE. Edit canonical lessons, then run npm run generate:books.
% canonical-source-hash: fnv1a64:76ce493c96b726b4
% canonical-lessons: TE-C182-nissabdam, TE-C182-mannikaina, TE-C182-nakili, TE-C182-sada, TE-C182-adanapu

\chapter{Silence, Durable, Fake, Plain, Extra}
\label{ch:te-silence-durable-fake-plain-extra}
\hlchaptermodality{182}

\begin{chapteropening}
\textbf{By the end of this chapter:} \emph{I can say silence, durable, fake, plain and extra.}

Complete the last lesson of chapter 182: I can say silence, durable, fake, plain and extra.
\end{chapteropening}

\section[\texorpdfstring{\te{నిశ్శబ్దం} (niśśabdaṁ)}{నిశ్శబ్దం (niśśabdaṁ)}]{\te{నిశ్శబ్దం} (niśśabdaṁ) — silence; quiet}

\label{lesson:TE-C182-nissabdam}



\begin{quote}
Before the new one: say the Telugu for wonderful, then the Telugu for ordinary.
\end{quote}

\subsection*{You'll want to know: \te{నిశ్శబ్దం}}
\textbf{\te{నిశ్శబ్దం}} — \emph{niśśabdaṁ} — \textquotedblleft{}silence; quiet\textquotedblright{}.

\begin{grammarlens}[title={What you've built}]
One more word for this chapter.
\end{grammarlens}

\subsection*{Your turn}
\begin{itemize}
\raggedright
  \item \emph{Say it:} \emph{niśśabdaṁ}
  \item \emph{Say it:} \emph{niśśabdaṁ}, once more
  \item \emph{Say it:} say \emph{niśśabdaṁ}
  \item \emph{Recall:} say the Telugu for wonderful, then the Telugu for ordinary, then say \emph{niśśabdaṁ} again
\end{itemize}

\begin{tcolorbox}[breakable,colback=teal!4,colframe=teal!35!black,title={Before you move on}]
What does \te{నిశ్శబ్దం} mean? (\textquotedblleft{}Silence; quiet\textquotedblright{}.) Say it once more.
\end{tcolorbox}
\section[\texorpdfstring{\te{మన్నికైన} (mannikaina)}{మన్నికైన (mannikaina)}]{\te{మన్నికైన} (mannikaina) — durable, long-lasting}

\label{lesson:TE-C182-mannikaina}



\begin{quote}
Before the new one: say the Telugu for ordinary, then the Telugu for silence.
\end{quote}

\subsection*{You'll want to know: \te{మన్నికైన}}
\textbf{\te{మన్నికైన}} — \emph{mannikaina} — \textquotedblleft{}durable, long-lasting\textquotedblright{}.

\begin{grammarlens}[title={What you've built}]
One more word for this chapter.
\end{grammarlens}

\subsection*{Your turn}
\begin{itemize}
\raggedright
  \item \emph{Say it:} \emph{mannikaina}
  \item \emph{Say it:} \emph{mannikaina}, once more
  \item \emph{Say it:} say \emph{mannikaina}
  \item \emph{Recall:} say the Telugu for ordinary, then the Telugu for silence, then say \emph{mannikaina} again
\end{itemize}

\begin{tcolorbox}[breakable,colback=teal!4,colframe=teal!35!black,title={Before you move on}]
What does \te{మన్నికైన} mean? (\textquotedblleft{}Durable, long-lasting\textquotedblright{}.) Say it once more.
\end{tcolorbox}
\section[\texorpdfstring{\te{నకిలీ} (nakilī)}{నకిలీ (nakilī)}]{\te{నకిలీ} (nakilī) — fake, counterfeit}

\label{lesson:TE-C182-nakili}



\begin{quote}
Before the new one: say the Telugu for silence, then the Telugu for durable.
\end{quote}

\subsection*{You'll want to know: \te{నకిలీ}}
\textbf{\te{నకిలీ}} — \emph{nakilī} — \textquotedblleft{}fake, counterfeit\textquotedblright{}.

\begin{grammarlens}[title={What you've built}]
One more word for this chapter.
\end{grammarlens}

\subsection*{Your turn}
\begin{itemize}
\raggedright
  \item \emph{Say it:} \emph{nakilī}
  \item \emph{Say it:} \emph{nakilī}, once more
  \item \emph{Say it:} say \emph{nakilī}
  \item \emph{Recall:} say the Telugu for silence, then the Telugu for durable, then say \emph{nakilī} again
\end{itemize}

\begin{tcolorbox}[breakable,colback=teal!4,colframe=teal!35!black,title={Before you move on}]
What does \te{నకిలీ} mean? (\textquotedblleft{}Fake, counterfeit\textquotedblright{}.) Say it once more.
\end{tcolorbox}
\section[\texorpdfstring{\te{సాదా} (sādā)}{సాదా (sādā)}]{\te{సాదా} (sādā) — plain, simple}

\label{lesson:TE-C182-sada}



\begin{quote}
Before the new one: say the Telugu for durable, then the Telugu for fake.
\end{quote}

\subsection*{You'll want to know: \te{సాదా}}
\textbf{\te{సాదా}} — \emph{sādā} — \textquotedblleft{}plain, simple\textquotedblright{}.

\begin{grammarlens}[title={What you've built}]
One more word for this chapter.
\end{grammarlens}

\subsection*{Your turn}
\begin{itemize}
\raggedright
  \item \emph{Say it:} \emph{sādā}
  \item \emph{Say it:} \emph{sādā}, once more
  \item \emph{Say it:} say \emph{sādā}
  \item \emph{Recall:} say the Telugu for durable, then the Telugu for fake, then say \emph{sādā} again
\end{itemize}

\begin{tcolorbox}[breakable,colback=teal!4,colframe=teal!35!black,title={Before you move on}]
What does \te{సాదా} mean? (\textquotedblleft{}Plain, simple\textquotedblright{}.) Say it once more.
\end{tcolorbox}
\section[\texorpdfstring{\te{అదనపు} (adanapu)}{అదనపు (adanapu)}]{\te{అదనపు} (adanapu) — extra, additional}

\label{lesson:TE-C182-adanapu}



\begin{quote}
Before the new one: say the Telugu for fake, then the Telugu for plain.
\end{quote}

\subsection*{You'll want to know: \te{అదనపు}}
\textbf{\te{అదనపు}} — \emph{adanapu} — \textquotedblleft{}extra, additional\textquotedblright{}.

\begin{grammarlens}[title={What you've built}]
One more word for this chapter.
\end{grammarlens}

\subsection*{Your turn}
\begin{itemize}
\raggedright
  \item \emph{Say it:} \emph{adanapu}
  \item \emph{Say it:} \emph{adanapu}, once more
  \item \emph{Say it:} say \emph{adanapu}
  \item \emph{Recall:} say the Telugu for fake, then the Telugu for plain, then say \emph{adanapu} again
\end{itemize}

\begin{tcolorbox}[breakable,colback=teal!4,colframe=teal!35!black,title={Before you move on}]
What does \te{అదనపు} mean? (\textquotedblleft{}Extra, additional\textquotedblright{}.) Say it once more.
\end{tcolorbox}
